% Generated draft: requires mathematical and source review.
\documentclass[11pt]{article}
\usepackage[margin=25mm]{geometry}
\usepackage{fontspec}
\setmainfont{Latin Modern Roman}
\usepackage{amsmath,amssymb,mathtools,unicode-math}
\setmathfont{Latin Modern Math}
\usepackage{enumitem,needspace,xurl,hyperref}
\setlength{\parindent}{0pt}
\setlength{\parskip}{5pt}
\setlength{\emergencystretch}{3em}
\newcommand{\N}{\mathbb{N}}
\newcommand{\Z}{\mathbb{Z}}
\newcommand{\Q}{\mathbb{Q}}
\newcommand{\R}{\mathbb{R}}
\newcommand{\C}{\mathbb{C}}
\newcommand{\F}{\mathbb{F}}
\newcommand{\A}{\mathbb{A}}
\newcommand{\PP}{\mathbb{P}}
\newcommand{\E}{\mathbb{E}}
\newcommand{\eps}{\varepsilon}
\newcommand{\dd}{\,\mathrm{d}}
\begin{document}
\section*{UnsolvedMath ID: 20003293}
\textbf{Source problem code:} catalogue entry
\subsection{Definitions and mathematical statement}
This record concatenates the end of the status discussion of Conjecture 21, the setup of section 4.5, and the introduction to Question 22; it states no independent question. Its first sentence reports that Conjecture 21 holds for Lebesgue-almost every $\alpha$ for any positive $\nu$ [136], in particular $\nu=2$, giving another proof of [164]. Neither that conjecture nor the meaning of $\nu$ is supplied in this extract.

The new setup is a compact Riemannian manifold $M$ with a topologically transitive $C^1$ Anosov flow $\varphi^t$ and a continuous invariant splitting
\[
 TM=E^0\oplus E^s\oplus E^u,
\]
where $E^0$ is the one-dimensional flow direction and for some $C,\lambda>0$ and every $t\ge0$,
\[
 \|D\varphi^t v\|\le Ce^{-\lambda t}\|v\|\quad(v\in E^s),\qquad
 \|D\varphi^t w\|\ge C^{-1}e^{\lambda t}\|w\|\quad(w\in E^u).
\]
Its strong stable leaves are
\[
 W^{ss}(x)=\{y:d(\varphi^t x,\varphi^t y)\to0\text{ as }t\to+\infty\}.
\]
The source additionally assumes $E^s=E^s_+\oplus E^s_-$ continuously and invariantly, with $\mu_+>\mu_->\lambda$ and
\[
 \|D\varphi^t v\|\le Ce^{-\mu_+t}\|v\|\quad(v\in E^s_+),\qquad
 \|D\varphi^t w\|\ge C^{-1}e^{-\mu_-t}\|w\|\quad(w\in E^s_-).
\]
The last domain is an explicit consistent reading: the supplied aid reports that the official source itself prints $E^u_-$ there, although the splitting is of $E^s$. The source says $E^s_+$ integrates to a fast stable foliation, gives the geodesic flow of $\mathbb C H^2$ as a basic example, and introduces a nonlinear analogue of Raghunathan's unipotent-measure question, without including Question 22 itself.

The source also attributes unique holonomy-invariant transverse measure for the strong stable foliation to [28]. Exact hypotheses and normalization, and any distinction between general transitivity and further assumptions, are not supplied. This historical assertion is recorded as such rather than used as an unconditional theorem below.

The aid's concrete model uses
\[
 B=\begin{pmatrix}2&1\\1&1\end{pmatrix},\quad
 \beta=\frac{3+\sqrt5}{2},\quad Bv=\beta^{-1}v,\quad v\ne0,
\]
and the unit-roof suspension
\[
 M_B=(\mathbb T^2\times\mathbb T^2\times\R)/
 \bigl((x,y,s)\sim(B^2x,By,s-1)\bigr),\qquad
 \varphi^t[x,y,s]=[x,y,s+t].
\]
There are two stable rates, $2\log\beta$ and $\log\beta$. The quotient-compatible fast leaf action is
\[
 U_r[x,y,s]=[x+r\beta^{2s}v,y,s],\quad r\in\R.
\]
The reported classification says that all probabilities invariant under both $\varphi^t$ and every $U_r$ are precisely the measures induced on the suspension by
\[
 m_{\mathbb T^2}\otimes\eta\otimes\dd s\quad(0\le s<1),
 \qquad B_*\eta=\eta,
\]
where $\eta$ ranges over all $B$-invariant probabilities on the slow torus. Joint ergodicity is equivalent to $B$-ergodicity of $\eta$. Adding invariance under the compatible slow action
\[
 V_r[x,y,s]=[x,y+r\beta^s v,s]
\]
forces $\eta=m_{\mathbb T^2}$. This is a model classification, not a solution of an unstated general Question 22.
\subsection{Short English statement}
The extract provides Anosov fast/slow stable-foliation context and stops before the next question. In its supplied suspension model, flow and fast-leaf invariance leave an arbitrary invariant measure on the slow torus; adding slow-leaf invariance forces Haar measure.
\subsection{Sources}
\begin{itemize}
\item American Institute of Mathematics, \emph{Emerging applications of measure rigidity}, boundary between the Conjecture 21 discussion and section 4.5 before Question 22, \url{https://aimath.org/WWN/measrigid/measrigid.pdf}. Supplied bibliography labels: [136], [164], and [28]; full entries are absent. The source invokes Raghunathan's question and the geodesic flow of $\mathbb C H^2$.
\item Local UnsolvedMath registry, permanent ID $20003293$, source code \texttt{catalogue entry}, canonical location \texttt{aim-other-notes.json}, \texttt{notes[99]}, dated 14 August 2026. Statement status is reconstructed and unverified. Research attempt 1: context only, ``not a problem'' at run, proved rigor, low novelty confidence; machine-generated and not peer reviewed.
\end{itemize}
\end{document}
